\documentclass[11pt,a4paper]{article}
\usepackage[utf8]{inputenc}
\usepackage[T1]{fontenc}
\usepackage[spanish]{babel}
\usepackage{amsmath,amssymb,mathtools}
\usepackage{geometry}
\usepackage{enumitem}
\usepackage{lmodern}
\usepackage{microtype}
\geometry{margin=2.4cm}
\setlength{\parindent}{0pt}
\setlength{\parskip}{0.45em}
\newcommand{\R}{\mathbb{R}}
\newcommand{\tr}{\operatorname{tr}}
\newcommand{\diag}{\operatorname{diag}}
\begin{document}

\begin{center}
{\LARGE\bfseries Álgebra Lineal}\\[0.4em]
{\Large Tema 7: Cónicas y Cuádricas --- Soluciones}
\end{center}

\section{Ejercicio 1}
\subsection{Solución}

La matriz asociada a la parte cuadrática es:

\[
A=
\begin{pmatrix}
8&2\\
2&5
\end{pmatrix}.
\]

Su polinomio característico es:

\[
\det(A-\lambda I)
=
(8-\lambda)(5-\lambda)-4.
\]

Por tanto:

\[
\lambda_1=9,
\qquad
\lambda_2=4.
\]

Tomamos los autovectores:

\[
v_1=(2,1),
\qquad
v_2=(1,-2).
\]

Normalizando:

\[
P=
\frac1{\sqrt5}
\begin{pmatrix}
2&1\\
1&-2
\end{pmatrix}.
\]

Entonces:

\[
P^tAP=
\begin{pmatrix}
9&0\\
0&4
\end{pmatrix}.
\]

Con:

\[
\begin{pmatrix}
x\\y
\end{pmatrix}
=
P
\begin{pmatrix}
x'\\y'
\end{pmatrix},
\]

la ecuación queda:

\[
9x'^2+4y'^2
+
\frac{12}{\sqrt5}x'
-
\frac{24}{\sqrt5}y'
+4=0.
\]

Completando cuadrados:

\[
x''=
x'+\frac{2}{3\sqrt5},
\qquad
y''=
y'-\frac{3}{\sqrt5}.
\]

Se obtiene:

\[
\boxed{
9x''^2+4y''^2-4=0.
}
\]

Equivalentemente:

\[
\boxed{
\frac{x''^2}{4/9}
+
\frac{y''^2}{1}
=
1.
}
\]

Por tanto:

\[
\boxed{\text{Elipse real}.}
\]

\section{Ejercicio 2}
\subsection{Solución}

La matriz asociada es:

\[
A=
\begin{pmatrix}
0&\frac12\\
\frac12&0
\end{pmatrix}.
\]

Los autovalores son:

\[
\lambda_1=\frac12,
\qquad
\lambda_2=-\frac12.
\]

Tomamos:

\[
v_1=(1,1),
\qquad
v_2=(1,-1).
\]

Así:

\[
P=
\frac1{\sqrt2}
\begin{pmatrix}
1&1\\
1&-1
\end{pmatrix}.
\]

Después de la rotación:

\[
\frac12x'^2
-
\frac12y'^2
-
\sqrt2\,x'
-1=0.
\]

Completando cuadrados:

\[
x''=x'-\sqrt2,
\qquad
y''=y'.
\]

Por tanto:

\[
\boxed{
\frac12x''^2
-
\frac12y''^2
-2=0.
}
\]

Equivalente a:

\[
\boxed{
\frac{x''^2}{4}
-
\frac{y''^2}{4}
=
1.
}
\]

Por tanto:

\[
\boxed{\text{Hipérbola}.}
\]

\section{Ejercicio 3}
\subsection{Solución}

La matriz asociada es:

\[
A=
\begin{pmatrix}
3&1\\
1&3
\end{pmatrix}.
\]

Los autovalores son:

\[
\lambda_1=4,
\qquad
\lambda_2=2.
\]

Tomamos:

\[
v_1=(1,1),
\qquad
v_2=(1,-1).
\]

Por tanto:

\[
P=
\frac1{\sqrt2}
\begin{pmatrix}
1&1\\
1&-1
\end{pmatrix}.
\]

La ecuación rotada es:

\[
4x'^2+2y'^2-8\sqrt2\,x'+6=0.
\]

Completando cuadrados:

\[
x''=x'-\sqrt2,
\qquad
y''=y'.
\]

Así:

\[
\boxed{
4x''^2+2y''^2-2=0.
}
\]

Equivalentemente:

\[
\boxed{
2x''^2+y''^2=1.
}
\]

Por tanto:

\[
\boxed{\text{Elipse real}.}
\]

\section{Ejercicio 4}
\subsection{Solución}

Como no aparece término $xy$, no es necesaria una rotación.

Completamos cuadrados:

\[
x^2+6x
=
(x+3)^2-9.
\]

Entonces:

\[
(x+3)^2-9+5y+14=0.
\]

\[
(x+3)^2+5(y+1)=0.
\]

Tomando:

\[
x''=x+3,
\qquad
y''=y+1,
\]

se obtiene:

\[
\boxed{
x''^2=-5y''.
}
\]

Por tanto:

\[
\boxed{\text{Parábola}.}
\]

\section{Ejercicio 5}
\subsection{Solución}

La matriz de la parte cuadrática es:

\[
A=
\begin{pmatrix}
m&1\\
1&m
\end{pmatrix}
\]

y la matriz homogénea:

\[
\widetilde A=
\begin{pmatrix}
m&0&1\\
0&m&1\\
1&1&m
\end{pmatrix}.
\]

Los invariantes son:

\[
I_1=2m,
\]

\[
I_2=m^2-1,
\]

\[
I_3
=
\det\widetilde A
=
m(m^2-2).
\]

Si:

\[
m\neq0,\pm\sqrt2,
\]

la cónica es no degenerada.

Para:

\[
|m|<1,
\qquad
m\neq0,
\]

se tiene:

\[
I_2<0.
\]

Por tanto:

\[
\boxed{\text{Hipérbola}.}
\]

Para:

\[
m=\pm1,
\]

\[
I_2=0,
\qquad
I_3\neq0.
\]

Por tanto:

\[
\boxed{\text{Parábola}.}
\]

Para:

\[
1<|m|<\sqrt2,
\]

\[
I_2>0,
\qquad
I_1I_3<0.
\]

Por tanto:

\[
\boxed{\text{Elipse real}.}
\]

Para:

\[
|m|>\sqrt2,
\]

\[
I_2>0,
\qquad
I_1I_3>0.
\]

Por tanto:

\[
\boxed{\text{Elipse imaginaria}.}
\]

Si:

\[
m=0,
\]

\[
I_3=0,
\qquad
I_2<0.
\]

Por tanto:

\[
\boxed{\text{Dos rectas secantes}.}
\]

Si:

\[
m=\pm\sqrt2,
\]

\[
I_3=0,
\qquad
I_2>0.
\]

Por tanto:

\[
\boxed{\text{Un punto}.}
\]

En conjunto:

\[
\boxed{
\begin{array}{c|c}
m & \text{cónica}\\
\hline
|m|>\sqrt2 & \text{elipse imaginaria}\\
1<|m|<\sqrt2 & \text{elipse real}\\
m=\pm1 & \text{parábola}\\
|m|<1,\ m\neq0 & \text{hipérbola}\\
m=0 & \text{dos rectas secantes}\\
m=\pm\sqrt2 & \text{un punto}
\end{array}
}
\]

\section{Ejercicio 6}
\subsection{Solución}

La matriz asociada es:

\[
A=
\begin{pmatrix}
3&1&0\\
1&3&0\\
0&0&-2
\end{pmatrix}.
\]

Sus autovalores son:

\[
-2,\qquad2,\qquad4.
\]

Tomamos los autovectores:

\[
v_1=(0,0,1),
\]

\[
v_2=(1,-1,0),
\]

\[
v_3=(1,1,0).
\]

Una base ortonormal es:

\[
e_1=(0,0,1),
\]

\[
e_2=
\frac1{\sqrt2}(1,-1,0),
\]

\[
e_3=
\frac1{\sqrt2}(1,1,0).
\]

Por tanto:

\[
P=
\begin{pmatrix}
0&\frac1{\sqrt2}&\frac1{\sqrt2}\\
0&-\frac1{\sqrt2}&\frac1{\sqrt2}\\
1&0&0
\end{pmatrix}.
\]

Después de la rotación:

\[
-2x'^2+2y'^2+4z'^2+8x'+4y'-8=0.
\]

Completando cuadrados:

\[
x''=x'-2,
\qquad
y''=y'+1,
\qquad
z''=z'.
\]

Se obtiene:

\[
\boxed{
-2x''^2+2y''^2+4z''^2-2=0.
}
\]

Equivalentemente:

\[
\boxed{
-x''^2+y''^2+2z''^2=1.
}
\]

Por tanto:

\[
\boxed{\text{Hiperboloide de una hoja}.}
\]

\section{Ejercicio 7}
\subsection{Solución}

La matriz asociada es:

\[
A=
\begin{pmatrix}
1&0&-2\\
0&1&0\\
-2&0&1
\end{pmatrix}.
\]

Sus autovalores son:

\[
1,\qquad3,\qquad-1.
\]

Tomamos:

\[
v_1=(0,1,0),
\]

\[
v_2=(1,0,-1),
\]

\[
v_3=(1,0,1).
\]

Una base ortonormal es:

\[
e_1=(0,1,0),
\]

\[
e_2=
\frac1{\sqrt2}(1,0,-1),
\]

\[
e_3=
\frac1{\sqrt2}(1,0,1).
\]

Después de la rotación:

\[
x'^2+3y'^2-z'^2-4x'+2=0.
\]

Completando cuadrados:

\[
x''=x'-2,
\qquad
y''=y',
\qquad
z''=z'.
\]

Se obtiene:

\[
\boxed{
x''^2+3y''^2-z''^2-2=0.
}
\]

Equivalentemente:

\[
\boxed{
\frac{x''^2}{2}
+
\frac{3y''^2}{2}
-
\frac{z''^2}{2}
=
1.
}
\]

Por tanto:

\[
\boxed{\text{Hiperboloide de una hoja}.}
\]

\section{Ejercicio 8}
\subsection{Solución}

La matriz asociada es:

\[
A=
\begin{pmatrix}
0&0&2\\
0&1&0\\
2&0&0
\end{pmatrix}.
\]

Sus autovalores son:

\[
-2,\qquad2,\qquad1.
\]

Tomamos:

\[
v_1=(1,0,-1),
\]

\[
v_2=(1,0,1),
\]

\[
v_3=(0,1,0).
\]

Una base ortonormal es:

\[
e_1=
\frac1{\sqrt2}(1,0,-1),
\]

\[
e_2=
\frac1{\sqrt2}(1,0,1),
\]

\[
e_3=(0,1,0).
\]

Por tanto:

\[
P=
\begin{pmatrix}
\frac1{\sqrt2}&\frac1{\sqrt2}&0\\
0&0&1\\
-\frac1{\sqrt2}&\frac1{\sqrt2}&0
\end{pmatrix}.
\]

No hay términos lineales, por lo que la rotación proporciona directamente:

\[
\boxed{
-2x'^2+2y'^2+z'^2+1=0.
}
\]

Equivalentemente:

\[
\boxed{
2x'^2-2y'^2-z'^2=1.
}
\]

Por tanto:

\[
\boxed{\text{Hiperboloide de dos hojas}.}
\]

\end{document}
